% BANKS-SOURCE: pdf=72 print=62 kind=chapter id=chapter06
\BanksChapter{6}{Massive quantum electrodynamics}

% BANKS-SOURCE: pdf=72 print=62 kind=prose
In this chapter we will do some perturbative calculations in a theory called \emph{massive
QED}. It is quantum electrodynamics with the photon replaced by a massive vector
field. Actually, it is the simplest example of a theory with a Higgs mechanism, since we
have seen that the massive vector can be written as a gauge theory by introducing an
additional scalar degree of freedom. Indeed, we can write a theory of electrodynamics
of spinor and scalar fields in terms of the Lagrangian

% BANKS-SOURCE: pdf=72 print=62 kind=equation id=6.1
\begin{equation}
  g^2\mathcal{L}=-|D_\mu\phi|^2-\frac14F_{\mu\nu}F^{\mu\nu}
  +\bar\psi(-\ii\gamma^\mu D_\mu-m)\psi-\kappa(|\phi|^2-v^2)^2.
  \label{eq:6.1}
\end{equation}
If we write $\phi=\rho\ee^{\ii\theta}$, and treat $\theta$ as the Stueckelberg field in the gauge formalism for the
massive vector field, then we can eliminate $\theta$ in favor of a massive vector field $B_\mu$ with a
mass that depends on the field $\rho$. $\rho$ is called the Higgs field. In the formal limit $\kappa\to\infty$,
fluctuations around $\rho=v$ are infinitely costly, and we get the theory called massive
QED. It turns out to be a perfectly good quantum field theory (it is renormalizable in
perturbation theory), in the sense that, if $g^2$ is small and we introduce an ultraviolet
cut-off $\Lambda$ to make sure the theory is well defined (see the chapter on renormalization),
then as long as $\Lambda<\ee^{c/g^2}E$, where $c$ is a number of order 1 that we will learn to compute
later, the predictions of the theory at energy scales $E\ll\Lambda$ are insensitive to the value
of $\Lambda$ and to the precise way in which we implement the cut-off.

We introduce massive QED for several reasons. First, it allows us to do computa-
tions in QED without worrying about gauge invariance. This is because the massive
theory has no gauge invariance, and can be quantized in a straightforward manner. The
Heisenberg equations of motion and canonical commutation relations lead, in a famil-
iar fashion, to the path-integral formula for the generating functional of Euclidean
Green functions

% BANKS-SOURCE: pdf=72 print=62 kind=equation id=6.2
\begin{equation}
  Z[J^\mu,\eta,\bar\eta]
  =\frac{\int[dB_\mu\,d\psi\,d\bar\psi]\,
  \ee^{-S+\int\dd^D x(B_\mu J^\mu+\bar\eta\psi+\bar\psi\eta)}}
  {\int[dB_\mu\,d\psi\,d\bar\psi]\,\ee^{-S}}.
  \label{eq:6.2}
\end{equation}
This leads to Euclidean Feynman rules with the propagators for $B_\mu$:

% BANKS-SOURCE: pdf=72 print=62 kind=equation id=6.3
\begin{equation}
  \langle B_\mu(p_E)B_\nu(-p_E)\rangle
  =\frac{\delta_{\mu\nu}+p_{E\mu}p_{E\nu}/\mu^2}{p_E^2+\mu^2}.
  \label{eq:6.3}
\end{equation}
The Euclidean fermion propagator is
\begin{equation*}
  \frac{1}{\slashed p_E+\ii m}.
\end{equation*}

% BANKS-SOURCE: pdf=73 print=63 kind=prose
As usual we have written the propagators of fields that are rescaled to make the
quadratic terms in the Lagrangian $e$-independent. The gauge boson mass is $\mu=ev$.
The interaction vertex corresponds to the following amputated\footnote[1]{This adjective implies, as usual, that ``the legs have been cut-off.''} Feynman diagram:

% BANKS-SOURCE: pdf=73 print=63 kind=figure id=unnumbered-main-ch6-fermion-gauge-vertex
\begin{center}
% Native TikZ for the unnumbered Chapter 6 fermion gauge vertex.
\begin{tikzpicture}[x=1cm,y=1cm,>=Stealth,baseline=(current bounding box.center)]
  \tikzset{BanksFermion/.style={line width=.55pt,postaction={decorate},
      decoration={markings,mark=at position .55 with {\arrow{Stealth[length=4pt]}}}}}
  % fermion line: in from upper left, out to upper right
  \draw[BanksFermion] (-0.62,0.78) -- (0,0);
  \draw[BanksFermion] (0,0) -- (0.62,0.78);
  % gauge boson leg
  \draw[decorate,decoration={snake,amplitude=1.8pt,segment length=4.5pt},line width=.55pt]
    (0,0) -- (0,-0.70);
  \node[anchor=west] at (3.6,0.30) {$\ii\ee\gamma^\mu$};
  \node[anchor=west] at (5.0,0.30) {$-\ee\gamma^\mu$};
\end{tikzpicture}

\end{center}
\BanksSourceNote{Fermion gauge vertex}

% BANKS-SOURCE: pdf=73 print=63 kind=prose
and has the value $-e\gamma^\mu$, which goes to $-\ii e\gamma^\mu$ in Lorentzian signature. It is a non-trivial,
but true, statement that, when we restrict attention to sources satisfying $\partial_\mu J^\mu=0$, this
generating functional has a finite limit as $\mu\to0$ and the vector boson becomes massless.
In this limit the particle states split into two different representations of the Poincaré
group, a massless scalar (the erstwhile longitudinal component of the massive spin-1
particle) and a massless particle with helicity $\pm1$. The restriction to conserved sources
decouples the massless scalar. A particular class of conserved currents, of the form
$J^\mu=\partial_\nu M^{\mu\nu}$, with $M^{\mu\nu}=-M^{\nu\mu}$, generates Green functions of the electromagnetic
field strength $F_{\mu\nu}$ in the $\mu\to0$ limit.

When we continue to Lorentzian signature and compute scattering amplitudes we see
several interesting effects in the massless limit. First of all the amplitudes for producing
longitudinal gauge bosons from initial states that had no longitudinal gauge bosons in
them goes to zero in this limit. Secondly, amplitudes involving scattering of fermions
diverge in perturbation theory. This divergence is primarily due to the fact that in the
$\mu\to0$ limit one can emit an arbitrary number of vector bosons with a finite cost in
energy. Indeed, in this limit the probability of emitting any finite number of bosons goes
to zero.\footnote[2]{To understand how a zero can look infinite in perturbation theory, consider the expression $\ee^{-e^2(\ln\mu)^2}$.} This is the \emph{infrared catastrophe} of Maxwell’s QED. As a catastrophe, it rates
pretty low on the Richter scale. It simply means that \emph{any} scattering of charged particles
is accompanied by a low-energy burst of classical \emph{bremsstrahlung} radiation, which more
or less follows the trajectories of the outgoing charged particles. If we restrict the energy
$\epsilon$ in this radiation, rather than the number of particles (which is always infinite in the
$\mu\to0$ limit) we get a finite answer. The advantage of massive QED is that we can
derive this finite answer by summing up finite amplitudes in perturbation theory and
then taking $\mu$ to zero, rather than formally resumming a series of infinite terms.

We do not have space in this text to actually do a bremsstrahlung calculation, and
will leave it to the problem set. Instead we will do two tree-level computations of
fermion annihilation processes, and a calculation through one loop level of low-energy
scattering in an external magnetic field. It will be convenient to add a second fermion
to the theory, with the same charge as the original one, and a mass $M\gg m$. We will call
the light fermions electrons and positrons, and the heavy ones muons and anti-muons,
in honor of famous characters who have appeared in the \emph{Real World}.\footnote[3]{A well-known piece of international cinema, of little artistic value, but enormous popular appeal.}
